\documentclass[11pt,a4paper]{article}
\usepackage[T1]{fontenc}
\usepackage[margin=2.2cm]{geometry}
\usepackage{graphicx}
\setlength{\parindent}{0pt}
\setlength{\parskip}{5pt}
\newcommand{\evidencia}[3]{%
  \textbf{#1. #2}\par
  \includegraphics[width=0.89\linewidth,height=3.8cm,keepaspectratio]{img/#1.png}\par
  #3\par\medskip
}
\begin{document}
\section*{AforoSeguro de práctica}
Práctica asistida con IA. La evidencia procede de procesos ejecutados en Linux.
Las imágenes reproducen las transcripciones de consola conservadas en
\texttt{docs/img/}.

\begin{tabular}{@{}p{\linewidth}@{}}
\evidencia{E0}{Conexión local}{El coordinador atendió en
\texttt{127.0.0.1} y recibió latidos de las tres puertas. La prueba
comprueba el uso de loopback en este entorno de práctica.}\strut\\
\evidencia{E1}{Operación normal}{Se aceptó una entrada en cada puerta y
una salida en P2. Las tres puertas observaron 2 de 10. La salida
previa con el contador en cero fue rechazada.}\strut\\
\evidencia{E2}{Límite y autenticación}{La undécima entrada fue rechazada.
El coordinador registró también el rechazo del token incorrecto y del
token ausente.}\strut
\end{tabular}

\newpage
\section*{Carga y fallos}
\begin{tabular}{@{}p{\linewidth}@{}}
\evidencia{E3}{Carga concurrente}{Con aforo 10, el modo normal permitió 10
entradas y rechazó 20. El modo inseguro permitió 30 y rechazó 0.
Las ocupaciones finales fueron 10 y 30, sin errores de transporte.}\strut\\
\evidencia{E4}{Caída y recuperación de P3}{P3 pasó a caída aproximadamente
seis segundos después de detener el proceso. El primer latido del proceso
reiniciado la devolvió a activa.}\strut\\
\evidencia{E5}{Coordinador detenido}{P1 informó que la entrada no se
había confirmado. Las puertas dependen de la respuesta del coordinador
para autorizar una entrada.}\strut
\end{tabular}

\newpage
\section*{Orden y análisis}
\begin{tabular}{@{}p{\linewidth}@{}}
\evidencia{E6}{Orden total}{La bitácora combinada contiene 1045 eventos. Las
solicitudes identificadas en la imagen se enviaron en las marcas 21 y 24.
El coordinador (P0) las recibió en 22 y 25, respectivamente. El orden
por Lamport mantiene ambos envíos antes de su recepción.}\strut
\end{tabular}

\textbf{1. Concurrencia.}
En E3, las solicitudes superaron la comprobación antes de que se aplicaran
los incrementos pendientes tras la pausa. El modo normal mantiene la
comprobación y el incremento en una sección crítica. Las operaciones
compiten por el mismo bloqueo y esperan su turno cuando hay contención.

\textbf{2. Consistencia y disponibilidad.}
En E5, P1 dejó de confirmar entradas. Si una partición produjera la misma
pérdida de comunicación, decidir localmente podría romper el límite global.
La decisión prioriza la consistencia del aforo ante esa pérdida de
comunicación. El público debe esperar y la administración conserva el
control central de las autorizaciones. Detener el coordinador en E5
demuestra la pérdida de servicio, mientras que CAP describe el compromiso
durante una partición.

\textbf{3. Coordinador único.}
Se podría replicar el estado con un líder elegido por mayoría y confirmar
cada operación tras replicar su posición en el registro. Las réplicas
tendrían que acordar el orden de entradas y salidas para evitar que dos
líderes concedan el último cupo. El prototipo guarda el contador en
memoria y lo pierde cuando se reinicia el coordinador.

\textbf{4. Latencia.}
La latencia de red puede ser significativa y variable. Los latidos
cada dos segundos dan margen para retrasos antes del umbral de seis
segundos. Con siete segundos de demora en una dirección, un latido
puede llegar tarde y la puerta puede aparecer caída temporalmente
aunque siga ejecutándose. También vencería la espera de respuesta del
cliente. Si luego se reciben latidos válidos, el coordinador vuelve a
marcar la puerta activa.
\end{document}
